\documentclass[conference]{IEEEtran}
\IEEEoverridecommandlockouts
\usepackage{cite}
\usepackage{amsmath,amssymb,amsfonts}
\usepackage{graphicx}
\usepackage{textcomp}
\usepackage{xcolor}
\usepackage{array}
\usepackage{booktabs}
\usepackage{listings}
\usepackage{hyperref}
\hypersetup{colorlinks=true, linkcolor=blue, citecolor=blue, urlcolor=blue}

\lstset{
  basicstyle=\ttfamily\footnotesize,
  breaklines=true,
  frame=single,
  rulecolor=\color{gray!40},
  backgroundcolor=\color{gray!5},
  keepspaces=true,
  columns=fullflexible,
}

\def\BibTeX{{\rm B\kern-.05em{\sc i\kern-.025em b}\kern-.08em
    T\kern-.1667em\lower.7ex\hbox{E}\kern-.125emX}}

\begin{document}

\title{Security Analysis of ML-Powered Android Apps:\\
A Static-Analysis Pipeline for On-Device ML\\
Vulnerability Detection}

\author{\IEEEauthorblockN{Bruna Vasconcelos}
\IEEEauthorblockA{\textit{School of Computing} \\
\textit{University of Utah} \\
Salt Lake City, UT, USA \\
bruna.gvasconcelos54@gmail.com}}

\maketitle

\begin{abstract}
On-device machine learning is increasingly embedded in mobile
applications, yet standard Android security tooling (FlowDroid,
TaintDroid, Soot) predates widespread on-device ML and carries no
awareness of ML-specific artifacts, SDKs, or data flows. We present a
static-analysis pipeline that augments APK analysis with ML-aware
detection: it identifies embedded model files via extension and
magic-byte matching, quantifies their exposure with Shannon entropy,
flags sensitive permissions mapped to ML data sources, and scans
decompiled Smali bytecode for ML-SDK imports and poisoning-surface
signatures. Detection outputs are scored against a Critical / High /
Medium / Info rubric. We apply the pipeline to a pilot corpus of five
F-Droid applications (274\,MB) and a synthetic positive-control APK
containing a public MobileNet~v1 model. The positive control is
correctly identified at High severity; none of the five real apps bundle
on-device ML. We treat the null result as a pilot observation rather
than a structural claim about F-Droid, since a five-app sample covers
$\approx$0.1\% of the catalog. The tool is architected to transfer
directly to a Play~Store corpus, where on-device ML is known to be more
prevalent; the main engineering change required is raising the Smali
scan cap and adding an obfuscation-resistant detector.
\end{abstract}

\begin{IEEEkeywords}
Android security, on-device machine learning, static analysis, data
poisoning, membership inference, Truth~Serum, APK analysis
\end{IEEEkeywords}

\section{Introduction}
The \emph{Truth~Serum} attack of Tram\`{e}r et al.~\cite{truthserum}
amplifies membership inference against a model by inserting mislabeled
near-duplicates of target samples into the training set. After
training, the poisoned target samples exhibit elevated loss, making
them trivially identifiable to an adversary. The attack presumes that
an adversary can contribute training data -- a presumption that maps
onto any on-device ML pipeline that collects user-supplied data for
retraining, personalization, or federated learning.

Existing mobile-ML security work measures orthogonal dimensions.
Xu~et~al.~\cite{xu2019} presented a first empirical study of deep
learning apps on smartphones, characterising \emph{presence} of ML
across 16{,}500 Play~Store apps. Sun~et~al.~\cite{modelxray} released
ModelXRay, a large-scale static analyser that quantifies
\emph{insufficient model protection}, showing that $41\%$ of ML apps
ship unprotected models and that $66\%$ of protected models are still
extractable via lightweight dynamic analysis. Standard Android taint
tools -- FlowDroid~\cite{flowdroid} and
TaintDroid~\cite{taintdroid} -- predate widespread on-device ML and
carry no ML-specific sources or sinks. None of these audits the
\emph{training-surface} precondition that enables Truth~Serum: whether
user-supplied data flows into on-device retraining or personalisation.
That gap is the project's target.

This paper makes three contributions:
\begin{itemize}
\item A static-analysis pipeline, implemented in under 500 lines of
  Python, that extends standard Android APK analysis with ML-aware
  detection.
\item A severity rubric for on-device ML vulnerabilities with
  reproducible evidence (per-file entropy, magic-byte identification,
  and SHA-256 APK hashes).
\item An empirical pilot on a five-app F-Droid corpus with a synthetic
  positive control, demonstrating that the detector fires as expected
  while documenting the dominant sources of false-negative risk
  (ProGuard/R8 obfuscation, Smali scan cap).
\end{itemize}

\section{Threat Model}
\label{sec:threat_model}

\noindent \textbf{Adversary.}~A malicious training-data contributor:
any party that can influence data ingested by an app's on-device
training, personalization, or federated-learning loop. The adversary
does not require code execution on the device; a poisoned input
delivered through the normal data-collection channel is sufficient
(the Truth~Serum setting).

\noindent \textbf{Defender.}~An app auditor with access to a publicly
distributed APK and no access to the running app, backend services, or
runtime instrumentation. Static analysis only.

\noindent \textbf{Assets at risk.}~Model weights (integrity and
confidentiality), the training-data pipeline (integrity), and user
personally identifiable information (PII) routed through inference.

\noindent \textbf{Out of scope.}~Runtime exploits, supply-chain attacks
on the build pipeline, backend-service compromises, and app repackaging
/ Frida-based dynamic analysis -- all deferred to future work under
the one-day pilot scope.

\section{Research Questions}
\label{sec:rq}

\noindent \emph{Primary (minimum viable deliverable):}
\begin{itemize}
\item \textbf{RQ1:} Do poisoning surfaces exist in distributed Android
  apps? Does user-submitted data flow into on-device retraining
  without sanitization or anomaly detection?
\item \textbf{RQ2:} Are model files exposed? Are \texttt{.tflite},
  \texttt{.onnx}, or similar artifacts bundled unencrypted inside APK
  assets with no integrity check on load?
\end{itemize}

\noindent \emph{Stretch:}
\begin{itemize}
\item \textbf{RQ3:} Is training data handled securely (world-readable
  directories, unsanitized input)?
\item \textbf{RQ4:} Are model update channels safe (TLS, signed
  updates)?
\end{itemize}

\section{Methodology}
\label{sec:method}

The analysis is static only. For each APK, the pipeline executes the
following phases:
\begin{enumerate}
\item \textbf{SHA-256 hash} the APK bytes, supporting reproducibility.
\item \textbf{Archive scan.} Unzip without decoding. Match file
  extensions in $\{\mathtt{.tflite}, \mathtt{.onnx}, \mathtt{.pb},
  \mathtt{.pt}, \mathtt{.pkl}, \mathtt{.h5}, \mathtt{.mlmodel},
  \mathtt{.pth}, \mathtt{.caffemodel}\}$. For each hit, read the first
  64\,KB and compute Shannon entropy plus magic-byte identification.
  An entropy above 7.9\,bits/byte is flagged as \emph{high-entropy}:
  consistent with encryption \emph{or} compression, which are
  indistinguishable from the bytes alone.
\item \textbf{Native-library scan.} Regex over archive entries for
  \texttt{libtensorflowlite*.so}, \texttt{libonnxruntime*.so},
  \texttt{libmlkit*.so}, \texttt{libmediapipe*.so}, and
  \texttt{libpytorch*.so}.
\item \textbf{\texttt{apktool} decode}, with a 300\,s per-APK timeout
  to bound the impact of adversarially constructed APKs.
\item \textbf{Manifest parse.}~\texttt{AndroidManifest.xml} is parsed
  with \texttt{defusedxml.ElementTree}, resisting XXE and
  entity-expansion attacks, rather than the naive regex approach often
  seen in similar pipelines. We extract the package name and flag a
  curated set of sensitive permissions mapped to ML data sources
  (CAMERA, RECORD\_AUDIO, ACCESS\_FINE\_LOCATION, READ\_CONTACTS,
  READ\_SMS, READ/WRITE\_EXTERNAL\_STORAGE, INTERNET).
\item \textbf{Smali pattern scan.} Capped at the first 5\,000 Smali
  files per APK, sorted alphabetically so the scan is deterministic.
  A warning is emitted when the cap truncates the scan -- silent
  truncation would mask false negatives. Two pattern families are
  applied: ML-SDK imports (\texttt{org/tensorflow/lite},
  \texttt{com/google/mlkit}, \texttt{ai/onnxruntime},
  \texttt{com/google/mediapipe}, \texttt{org/pytorch}) and
  poisoning-surface signatures (\texttt{modelPersonalizer} /
  \texttt{PersonalizationModel} / \texttt{OnDeviceCustomActions},
  federated-learning references, feedback-to-retraining phrases,
  \texttt{GradientTape} / \texttt{keras.fit}, and
  \texttt{differentialPrivacy} APIs).
\item \textbf{Severity scoring.} Findings are scored Critical / High /
  Medium / Info by the rubric defined in the proposal.
\end{enumerate}

The pipeline emits one JSON record per APK (including per-APK SHA-256
and both scanned and total Smali file counts) plus an aggregated CSV.

\subsection{Scope cuts for the one-day pilot}

The full methodology documented in the project proposal includes
FlowDroid-based taint analysis and Frida dynamic validation. Both are
deferred: FlowDroid's configuration cost outweighs its yield on a
one-day timeline, and its known reflection / dynamic-class-loading
blind spots would require manual review regardless. Frida requires
APK repackaging, estimated at approximately two days per batch. The
pilot corpus is reduced from 30+ apps to five real apps plus one
synthetic positive-control APK.

\section{Corpus and Ethics}
\label{sec:corpus}

The pilot corpus consists of five applications downloaded from
F-Droid (open-source, freely redistributable, no CFAA concern) plus
one synthetic positive-control APK assembled from a public
MobileNet~v1 \texttt{.tflite} file (Apache 2.0 license, Google Cloud
Storage). Table~\ref{tab:corpus} lists the corpus and SHA-256
prefixes for reproducibility.

\begin{table}[htbp]
\caption{Pilot Corpus with SHA-256 Prefixes}
\label{tab:corpus}
\centering
\footnotesize
\begin{tabular}{@{}lll@{}}
\toprule
\textbf{App} & \textbf{Size} & \textbf{SHA-256 (16)} \\
\midrule
F-Droid (\texttt{org.fdroid.fdroid}) & 12.4\,MB & \texttt{985f5181d48bb6ba} \\
NewPipe (\texttt{org.schabi.newpipe}) & 10.9\,MB & \texttt{dbc8a1bb7a3db16f} \\
Aves~Gallery libre & 55.9\,MB & \texttt{75697b19f2eb850f} \\
Nextcloud Files & 77.8\,MB & \texttt{e36e6ef4215cf003} \\
Fennec (Firefox fork) & 117.6\,MB & \texttt{251146a2b5f6d801} \\
\textbf{Positive control} (synthetic) & 15.7\,MB & \texttt{d9c71359fcce2d5f} \\
\bottomrule
\end{tabular}
\end{table}

All analysis is performed on freely redistributable APKs. No server
interaction, no live user data, no repackaging, and no PoC delivery
to real devices. Responsible disclosure is framed as a 90-day window
should any critical finding emerge in future Play~Store runs.

\section{Findings}
\label{sec:findings}

\subsection{Pilot observation on the F-Droid corpus}

None of the five F-Droid APKs bundle a model file, link an ML native
library, or reference an ML SDK in the Smali that was scanned. Every
app's top finding is an Info-severity \emph{``No on-device ML
artifacts detected.''}

We treat this as a \emph{pilot observation}, not a structural claim
about F-Droid as a whole. Five apps is approximately 0.1\% of the
F-Droid catalog, and the Smali scan is capped
(\S\ref{sec:findings_coverage}); absence of evidence at this scale is
not evidence of absence. The defensible reading is that F-Droid's
audit requirements \emph{likely} filter out proprietary on-device ML,
which should be confirmed on a larger corpus.

The actionable conclusion: F-Droid is not a productive corpus for
this research question. The pipeline should be applied to Play~Store
apps in high-risk categories (\S\ref{sec:conclusion}).

\subsection{Positive control: detection fires as expected}

The synthetic positive-control APK, containing
\texttt{assets/mobilenet\_v1.tflite} (16.9\,MB, Google MobileNet~v1),
is correctly identified by the pipeline:

\begin{itemize}
\item Model artifact detected at \texttt{assets/mobilenet\_v1.tflite}.
\item Entropy 7.43\,bits/byte, flagged as \emph{structured, not high-
  entropy} (as expected for a stock TFLite model).
\item Magic bytes matched; format identified as \texttt{tflite}.
\item Severity: \textbf{High} (``unencrypted model present;
  architecture and weights readable by anyone who can unzip the
  APK'').
\end{itemize}

This establishes that the null result on the F-Droid corpus reflects
the absence of ML artifacts in the scanned files, not a tool failure.

\subsection{Per-app severity and Smali coverage}
\label{sec:findings_coverage}

Table~\ref{tab:findings} reports per-app severity counts and Smali
scan coverage. The coverage column is the most important new datum:
the three largest apps had only 11--25\% of their Smali files scanned
under the default 5\,000-file cap. A missed ML-SDK import in the
untaken 75--89\% would be a silent false negative. Raising the cap is
a trivial configuration change for the Play~Store pivot and is listed
as the top item in \S\ref{sec:limitations}.

\begin{table}[htbp]
\caption{Per-App Severity and Smali Coverage}
\label{tab:findings}
\centering
\footnotesize
\begin{tabular}{@{}lrrrrr@{}}
\toprule
\textbf{App} & \textbf{Smali} & \textbf{Crit} & \textbf{High} & \textbf{Med} & \textbf{Info} \\
 & \textbf{cov.\%} & & & & \\
\midrule
aves & 75 & 0 & 0 & 0 & 1 \\
fdroid & 25 & 0 & 0 & 0 & 1 \\
fennec & 15 & 0 & 0 & 0 & 1 \\
newpipe & 58 & 0 & 0 & 0 & 1 \\
nextcloud & 11 & 0 & 0 & 0 & 1 \\
pos.\ ctrl. & n/a & 0 & \textbf{1} & 0 & 0 \\
\bottomrule
\end{tabular}
\end{table}

All-Info is the correct outcome when no ML artifacts are present; the
tool also surfaces sensitive-permission posture
(\S\ref{sec:findings_perms}) and records per-APK SHA-256 for
reproducibility.

\subsection{Sensitive-permission context}
\label{sec:findings_perms}

The pipeline reported sensitive-permission posture for every app even
where no ML artifacts were present. Fennec declared the most
sensitive permissions (6), followed by Nextcloud (5), F-Droid (4),
Aves (3), and NewPipe (2). These permissions are \emph{latent} risk:
if any of these apps added on-device ML in a future release, the data
pipeline would already have access to CAMERA, MIC, FINE\_LOCATION, or
INTERNET without re-requesting consent.

\section{Tool Limitations}
\label{sec:limitations}

\begin{enumerate}
\item \textbf{Smali scan cap.} Nextcloud's 46\,547-file Smali tree was
  scanned at only 11\% coverage; three of five apps fell below 25\%.
  The cap exists to bound pilot runtime and must be raised (or
  replaced with priority-ordered scanning) before any claim of full
  coverage.
\item \textbf{Obfuscation blindness.} ProGuard/R8 renames
  \texttt{org.tensorflow.lite.Interpreter} to, e.g., \texttt{a.b.c}
  before shipping. Pattern-based Smali detection is defeated. This is
  the dominant risk for the planned Play~Store pivot and requires
  either native-library-anchored heuristics (libraries are rarely
  obfuscated) or string-constant analysis of the resource table.
\item \textbf{Entropy is not a cryptographic signal.} High Shannon
  entropy is consistent with encryption \emph{or} compression; the
  pipeline labels the condition \texttt{high\_entropy} (not
  \texttt{likely\_encrypted}) and notes the ambiguity in every
  affected finding.
\item \textbf{No reflection / dynamic-class-loading analysis.} Smali
  pattern scan misses ML SDK use hidden behind \texttt{Class.forName}
  or DEX-at-runtime loaders -- the same blind spot FlowDroid
  documents.
\item \textbf{Static only.} Cannot observe runtime model downloads or
  federated-learning rounds.
\item \textbf{Magic-byte signature set is not exhaustive.} Proprietary
  or custom model formats may evade the set. Entropy is a backstop.
\end{enumerate}

\section{Reproducibility}

\begin{lstlisting}[language=bash]
cd cs6491/project
python3 pipeline/pipeline.py \
  --apk-dir apks/ \
  --out results
\end{lstlisting}

Outputs land in \texttt{results/summary.csv} plus one JSON per app.
Full decoded trees are written under \texttt{decompiled/}. Per-APK
SHA-256 is recorded in both JSON and CSV. Runtime was approximately
three minutes for the 274\,MB corpus on a single Kali Linux laptop
(Python 3.13, \texttt{apktool} 2.7.0, \texttt{jadx} 1.5.2).

\section{Conclusion and Next Steps}
\label{sec:conclusion}

The pilot validates two claims. First, the pipeline runs end-to-end
with tools that ship on standard Linux distributions
(\texttt{apktool}, \texttt{jadx}, Python plus \texttt{defusedxml}) --
no heroic toolchain is required. Second, the detector fires correctly
on a known-positive sample (MobileNet~v1) while returning a clean
null on a five-app F-Droid corpus.

The next run should target Play~Store apps in three high-risk
categories: keyboard apps (continuous typing data $\rightarrow$
retraining surface), translation apps (offline TFLite models), and
camera / filter apps (on-device classification on PII). Two pipeline
changes are required first: raise the Smali cap (trivial) and add an
obfuscation-resistant detector (moderate -- anchored on native-
library presence plus string-constant extraction).

Paired with a training-set defense that detects Truth~Serum poisons
(out of scope for this pilot), this auditor completes the loop between
lab-side defense research and field-side risk assessment by
identifying the real-world training pipelines where such poisons
could be planted.

\begin{thebibliography}{00}

\bibitem{truthserum}
F. Tram\`{e}r, R. Shokri, A. San Joaquin, H. Le, M. Jagielski,
S. Hong, and N. Carlini, ``Truth Serum: Poisoning Machine Learning
Models to Reveal Their Secrets,'' in \emph{Proc. ACM SIGSAC Conf.
Computer and Communications Security (CCS)}, 2022.

\bibitem{xu2019}
M. Xu, J. Liu, Y. Liu, F. X. Lin, Y. Liu, and X. Liu,
``A First Look at Deep Learning Apps on Smartphones,'' in
\emph{Proc. World Wide Web Conference (WWW)}, 2019.

\bibitem{modelxray}
Z. Sun, R. Sun, L. Lu, and A. Mislove, ``Mind Your Weight(s):
A Large-scale Study on Insufficient Machine Learning Model Protection
in Mobile Apps,'' in \emph{Proc. USENIX Security Symposium}, 2021.

\bibitem{flowdroid}
S. Arzt, S. Rasthofer, C. Fritz, E. Bodden, A. Bartel, J. Klein,
Y. Le Traon, D. Octeau, and P. McDaniel, ``FlowDroid: Precise
Context, Flow, Field, Object-sensitive and Lifecycle-aware Taint
Analysis for Android Apps,'' in \emph{Proc. ACM SIGPLAN Conf. Programming
Language Design and Implementation (PLDI)}, 2014.

\bibitem{taintdroid}
W. Enck, P. Gilbert, B. Chun, L. P. Cox, J. Jung, P. McDaniel, and
A. N. Sheth, ``TaintDroid: An Information-Flow Tracking System for
Realtime Privacy Monitoring on Smartphones,'' in \emph{Proc. USENIX
Symp.\ Operating Systems Design and Implementation (OSDI)}, 2010.

\bibitem{spectral}
B. Tran, J. Li, and A. Madry, ``Spectral Signatures in Backdoor
Attacks,'' in \emph{Advances in Neural Information Processing Systems
(NeurIPS)}, 2018.

\bibitem{shokri}
R. Shokri, M. Stronati, C. Song, and V. Shmatikov, ``Membership
Inference Attacks Against Machine Learning Models,'' in \emph{Proc.
IEEE Symp. Security and Privacy (S\&P)}, 2017.

\bibitem{mlkit}
Google, ``ML~Kit: On-Device Machine Learning for Mobile,''
2023.\ [Online]. Available: \url{https://developers.google.com/ml-kit}

\end{thebibliography}

\end{document}
